\chapter{Background}
\label{chap:background}

This chapter assembles the foundations the rest of the document rests on, in the order a reader
needs them. It begins with the object of the comparison: language models, the tokenisers that
turn text into the units they process, and the adaptation method that fits them to a task. It
then sets out the three conditions under which the models are compared: reduced weight precision,
inference on the processor of a \gls{sbc}, and decoding under a grammar. It closes
with the task those models are asked to perform, the extraction of an intent and its slots from a
spoken command, and with the instruments used to score them: the accuracy and efficiency metrics,
the paired statistical test, and the notion of dominance that organises a trade-off. Every section
states established material. What the literature has and has not measured under these conditions
is the subject of Chapter~\ref{chap:state-of-the-art}, and every design decision made in this work
is stated and justified in Chapter~\ref{chap:method}.

\section{Language models and the transformer decoder}
\label{sec:bg-lm}

A language model assigns a probability to a sequence of tokens $x_1, \dots, x_T$ by factorising
it into a product of conditional next-token distributions,
\begin{equation}
    p(x_1, \dots, x_T) = \prod_{t=1}^{T} p(x_t \mid x_1, \dots, x_{t-1}),
    \label{eq:bg-autoregressive}
\end{equation}
and an autoregressive model is one trained to estimate each factor from the tokens before it.
Generation runs the factorisation forward: given a prompt, the model produces a distribution over
the next token, one token is chosen from it, appended to the sequence, and the step repeats until
a stop token is produced or a length cap is reached. When the chosen token is always the most
probable one the procedure is \emph{greedy decoding}; for a fixed implementation, with one
sequence decoded at a time, its output for a given prompt is deterministic. Sampling at a temperature, or searching over several partial sequences at once,
trades that determinism for diversity. Two phases of unequal cost make up one generation. The
\emph{prefill} phase processes every prompt token in one pass; the \emph{decode} phase then
produces one output token per forward pass. The per-token cost of decoding, not the cost of
prefill, is what bounds how long a short structured output takes to produce.

\paragraph{The decoder-only transformer.} Every model compared in this work is a decoder-only
transformer, the architecture introduced for machine translation~\cite{vaswani2017} and since
adopted for language modelling by every model family compared here~\cite{qwen25,smollm2,llama32card,danube3}.
The architecture maps each input token to a vector through an embedding table, passes the
sequence of vectors through a stack of $L$ identical blocks, and projects the final vector at
each position onto the vocabulary to obtain a score per token, the \emph{logits}, which a softmax
turns into the next-token distribution of Equation~\ref{eq:bg-autoregressive}.
Each block has two sub-layers joined by residual connections and normalisation. The first is
multi-head self-attention: three linear projections produce a query, a key and a value for every
position, each query is scored against the keys of every earlier position, and the values are
combined under those scores and passed through a fourth, output projection. The four
projections are conventionally named \texttt{q\_proj}, \texttt{k\_proj}, \texttt{v\_proj} and
\texttt{o\_proj}. The second sub-layer is a position-wise feed-forward network, which in the
model family used here is gated: a \texttt{gate\_proj} and an \texttt{up\_proj} each expand the
vector; the gate branch is passed through a non-linearity and multiplied element-wise with the up
branch, and a \texttt{down\_proj} contracts the product back to the model width. These seven
projections hold most of a model's parameters, the embedding table holding the rest, and they are
the modules that the adaptation method of Section~\ref{sec:bg-lora} attaches to. A
causal mask restricts every position to attend to itself and earlier positions only, which is
what makes the stack an estimator of the conditional distributions in
Equation~\ref{eq:bg-autoregressive}.

\paragraph{The key--value cache.} During decoding, the attention of the newest position ranges
over the keys and values of every earlier one. Recomputing those for the whole sequence at every
step would make the cost of the $t$-th token grow with $t$. Instead the keys and values computed
at each step are cached, so that a decode step costs one forward pass over one token plus a read
of the cache. The cache grows with the context, and its maximum size is fixed by the context
window the runtime allocates. For the short prompts and outputs of a command parser the cache is
small, and its cost is dominated by the weights, not by the context.

\paragraph{Tokenisation.} A model does not read characters. A tokeniser segments text into pieces
drawn from a fixed vocabulary, typically of subwords learnt from a corpus by merging frequent
character sequences, and maps each piece to an integer identifier. Common words become one piece;
rare words, numbers and punctuation are split into several. A line of structured text such as a
\gls{json} object is therefore segmented into brace, quote, key, colon and value pieces, and the
number of tokens it occupies, not its length in characters, is what the decode phase pays for.
The segmentation is a property of the vocabulary and the merge rules together, and two
implementations of the same vocabulary are expected to produce the same identifiers for the same
string. A model is trained on one segmentation, so a serving implementation that segments the same
text differently presents the model with a sequence it never saw during training, with no error
raised. Chapter~\ref{chap:method} treats agreement between the training and serving tokenisers as
a property to be checked rather than assumed.

\paragraph{Instruction tuning and chat templates.} A \emph{base} model is trained only to
continue text. An \emph{instruct} or \emph{chat} model is further trained on conversations, so
that it responds to a request rather than continuing it. The conversation is rendered into one
token sequence by a \emph{chat template}: a fixed convention of special tokens and role markers
that delimits the user's turn, the assistant's reply and, where the template has one, a system
turn. The template is part of
the model's training distribution. Fine-tuning and serving must therefore render the same
template, and a difference in a role marker or a whitespace convention between the two changes
the prompt the model sees. Every model compared in this work is the instruct variant of its
family; how its template is rendered at fine-tuning and at serving, and how the two renderings are
checked against each other, is stated in Chapter~\ref{chap:method}.

\paragraph{Small language models.} The models compared here have between a third of a billion and
one and a quarter billion parameters. Such models are commonly called \glspl{slm} to distinguish
them from the \glspl{llm} of tens or hundreds of billions of parameters whose capabilities
motivated the field. The distinction is one of scale, not of architecture: an \gls{slm} is a
decoder-only transformer with fewer and narrower blocks, trained on a smaller or more carefully
curated corpus~\cite{smollm2}, and released in the same base and instruct variants. What scale
buys, at the cost of memory and time per token, is the question the comparison in this document
puts to four such models. Table~\ref{tab:bg-models} lists them with their parameter counts.

\begin{table}[htbp]
\centering
\caption[Small language models compared in this work]{The four small language models compared in
this work. Parameter counts are those of the released checkpoints, rounded to the nearest million;
the developer's report is cited for each model.}
\label{tab:bg-models}
\begin{tabular}{@{}llrl@{}}
\toprule
Model & Developer & Parameters & Source \\
\midrule
SmolLM2-360M-Instruct & Hugging Face & 362~M & \cite{smollm2} \\
Qwen2.5-0.5B-Instruct & Alibaba Qwen team & 494~M & \cite{qwen25} \\
H2O-Danube3-500M-Chat & H2O.ai & 514~M & \cite{danube3} \\
Llama-3.2-1B-Instruct & Meta & 1{,}236~M & \cite{llama32} \\
\bottomrule
\end{tabular}
\end{table}

% \noautobreak: the keep-section-whole rule would start this section on a new page
% and leave the rest of the page under Table~\ref{tab:bg-models} blank (layout pass, 25 Sep).
\noautobreak
\section{Fine-tuning and low-rank adaptation}
\label{sec:bg-lora}

A pretrained model is adapted to a task by \emph{supervised fine-tuning}: it is trained on pairs
of a prompt and the completion that the task expects, the loss being the negative logarithm of
the completion's probability under Equation~\ref{eq:bg-autoregressive}, the cross-entropy, taken
over the completion tokens. Computing the loss over
the prompt tokens as well would reward the model for reproducing input it is never asked to
generate; restricting the loss to the completion, called \emph{completion-only masking}, keeps
every gradient about whether the output that follows the prompt is the right one. Full
fine-tuning updates every parameter of the model. Its memory cost is several times the size of
the model itself, since the optimiser keeps state for every parameter it updates, and it produces
a complete copy of the weights per task.

\paragraph{Low-rank adaptation.} \Gls{lora} avoids both costs by freezing the pretrained weight
matrices and learning a low-rank update beside each targeted one~\cite{lora}. For a projection
$W \in \mathbb{R}^{d \times k}$, the adapted projection is
\begin{equation}
    W' = W + \frac{\alpha}{r}\, B A,
    \qquad B \in \mathbb{R}^{d \times r},\; A \in \mathbb{R}^{r \times k},\; r \ll \min(d, k),
    \label{eq:bg-lora}
\end{equation}
where only $A$ and $B$ are trained, $r$ is the adapter rank and $\alpha$ a scaling constant. The
trainable parameters number $r(d + k)$ per projection rather than $dk$, so an adapter for a
sub-billion model amounts to a few million parameters. After training, the product $BA$ can be
merged into $W$, so the adapted model has the same shape and the same inference cost as the
original. The adapter is attached to a chosen subset of the projections named in
Section~\ref{sec:bg-lm}; the choice of subset and rank are hyperparameters of the recipe.
QLoRA extends the method by quantising the frozen base weights to four bits during training,
keeping only the adapters and the optimiser state at higher precision, so that a model too large
for the training accelerator can still be adapted~\cite{qlora}. The same work found that adapters
on every linear projection, attention and feed-forward alike, were needed to match the accuracy of
full fine-tuning, whereas adapters on the query and value projections alone, the original
\gls{lora} default, fell short~\cite{qlora}.

\paragraph{Optimisation and checkpoint selection.} The adapters are trained with a
gradient-based optimiser; AdamW, which keeps a running estimate of each parameter's gradient mean
and variance and applies weight decay separately from the adaptive step, is the usual
choice~\cite{adamw}. The learning rate follows a schedule, commonly a short warm-up followed by a
cosine decay to zero. The model is evaluated on a held-out validation split at the end of every
epoch, and one checkpoint is kept. Validation loss is a per-token quantity and gives partial
credit to an output that is close in token terms yet wrong as a whole; an item-level task metric
does not. Which criterion selects the checkpoint is part of the recipe
(Chapter~\ref{chap:method}). The software stack
used for this stage in the present work, \texttt{transformers}~\cite{transformers},
\texttt{peft}~\cite{peft} and PyTorch~\cite{pytorch}, is named in Chapter~\ref{chap:method}
together with the recipe.

\section{Post-training quantisation and the GGUF formats}
\label{sec:bg-quant}

The weights of a trained model are stored as floating-point numbers, typically at 16~bits each.
\emph{Quantisation} replaces them by integers of fewer bits together with a small number of
floating-point scale factors, so that each weight $w$ is approximated as $w \approx s \cdot q$,
with $q$ an integer of $b$ bits and $s$ a scale shared by a block of weights. A block of 32
weights at 4~bits with one 16-bit scale occupies $32 \times 4 + 16 = 144$ bits, or 4.5~bits per
weight, against 512 for the same block at 16~bits. The approximation error is the rounding of $w/s$
to the nearest integer, and it is smaller for smaller blocks, which cost more scale factors, and
for more bits, which cost more storage. \emph{Post-training} quantisation applies this
transformation to a model after training, without further gradient updates; it is the route taken
when the training budget or the training data are no longer available at deployment time.
Quantisation-aware training, which simulates the rounding during training so that the weights
adapt to it, is the alternative, at the cost of a training run. Weight-only quantisation, which
leaves the activations at higher precision, is the common choice for inference at batch size one,
because at that batch size it is the reading of the weights, not the arithmetic on them, that
bounds the speed of decoding (Section~\ref{sec:bg-edge}), as a roofline analysis on a desktop
\gls{gpu} establishes~\cite{awq}.

\paragraph{Calibrated methods.} Rounding to the nearest integer treats every weight alike. The
methods introduced for \glspl{llm} use a small calibration set of inputs to reduce the error.
GPTQ quantises the weights of each layer one column at a time and adjusts the remaining columns to
compensate for the rounding error already committed, using second-order information about the
layer's inputs~\cite{gptq}. AWQ observes that a small fraction of weight channels, those that meet
large activations, account for most of the error, and scales them up before rounding so that
their relative precision is preserved~\cite{awq}. SpQR isolates the outlier weights that resist
low-bit representation and stores them separately at higher precision~\cite{spqr}. The cost of
any of these transformations is usually reported as the change in \emph{perplexity} on held-out
text, the exponential of the mean negative log-likelihood per token, or as the Kullback--Leibler
divergence of the quantised model's next-token distribution from the unquantised one. Both
measure how far the quantised model's next-token predictions have moved from the original's,
averaged over generic text; how that translates into accuracy on a specific downstream task is a
separate measurement.

\paragraph{GGUF and the llama.cpp formats.} \gls{gguf} is the file format of the
\texttt{llama.cpp} runtime (Section~\ref{sec:bg-edge}): one file holds every tensor of a model
together with the metadata a runtime needs to serve it, including the tokeniser vocabulary and the
chat template~\cite{llamacpp}. A model is converted to \gls{gguf} at 16~bits and then quantised by
the runtime's own tool to one of its formats. Two are used in this work. Q8\_0, a legacy format,
stores weights in blocks of 32 at eight bits with one scale per block, 8.5~bits per weight in
all. Q4\_K\_M belongs to the K-quant family, which stores weights in super-blocks of 256, each
divided into eight sub-blocks of 32 whose scales are themselves stored at reduced precision,
4.5~bits per weight at the four-bit level~\cite{llamacpp}. Its
target is four bits for most tensors. A fixed rule raises the output projection, and a subset of
the attention-value and feed-forward down-projection tensors chosen by layer position, to six
bits~\cite{llamacpp}. A K-quant also requires each tensor's rows to divide into whole
super-blocks: a tensor whose row length is not a multiple of 256 falls back to a legacy format,
five bits in place of four and eight in place of six~\cite{llamacpp}. The name Q4\_K\_M therefore
denotes an allocation recipe rather than a bit width, and in a model whose hidden dimension is
not a multiple of 256 most of the weights end up at five or eight bits. The recipe allocates bits
by tensor role, position and shape; unless an importance matrix computed from calibration text
is supplied, it uses no calibration data, which sets the plain form apart from the methods of the
previous paragraph. The runtime's documentation reports the cost of its formats as
perplexity or Kullback--Leibler divergence against the unquantised model~\cite{llamacpp}.

\section{Inference on a CPU-only single-board computer}
\label{sec:bg-edge}

A \gls{sbc} is a complete computer on one circuit board: a system-on-chip carrying the processor
cores, a memory controller and the peripheral interfaces, with the memory soldered beside it and
storage on a removable card. The Raspberry~Pi~5 used in this work carries a quad-core Arm
Cortex-A76 processor and LPDDR4X memory shared by all four cores~\cite{rpi5brief}, and no
accelerator that a language-model runtime uses in practice. Inference therefore runs on the
\gls{cpu} alone, through the vector instructions of the Arm architecture, and the four cores are
the only compute available to every process on the board: the speech recogniser, the language
model, the validator and the operating system share them. Running a language model on such a
board is the setting the phrase \emph{edge inference} denotes in this document: the model executes
on the device that uses its output, with no network round trip and no data-centre accelerator.

\paragraph{The llama.cpp runtime.} \texttt{llama.cpp} is a C/C++ inference engine for
transformer language models, built on the \gls{ggml} tensor library, that runs on a \gls{cpu}
without a \gls{gpu} driver or a Python interpreter~\cite{llamacpp}. It loads a \gls{gguf} file by
memory-mapping it, so that the weights are read from storage into memory on first use, and
executes the forward pass with kernels specialised for each quantisation format and for the
vector instruction set of the host. It exposes the number of threads it uses and can be pinned to
a subset of the cores by the operating system, which is how a share of the board is reserved for
it. It also implements the grammar-constrained sampling of Section~\ref{sec:bg-gbnf} and a server
that keeps a prompt cache, so that a prefix common to every request is prefilled once.

\paragraph{Memory-bound decoding.} A decode step at batch size one performs one
matrix--vector product per projection, and each such product reads the whole projection matrix
once to produce one vector. The arithmetic per weight is two operations, a multiply and an add,
while the weight itself must be fetched from memory. On a processor whose arithmetic throughput
exceeds its memory bandwidth by a wide margin, the step therefore finishes when the last weight
has been read, not when the last operation has been done, and the time per token is approximately
the size of the weights in bytes divided by the memory bandwidth in bytes per second. Two
consequences follow. First, under that bound a smaller model or a lower-precision format decodes
faster in proportion to the bytes it saves, which is the rationale for weight-only quantisation as
a means to decode throughput~\cite{awq}; how closely a given processor approaches the bound, given
the arithmetic its low-bit kernels add, is a matter of measurement. Second, adding threads beyond
what saturates the memory bus adds little, so the number of cores an inference process is given is a matter of how the board's
compute is shared rather than of how fast the model could run alone. Prefill is different: it
processes every prompt token in one matrix--matrix pass, reads each weight once for all of them,
and is bounded by arithmetic. For a short structured output following a short prompt, decode
dominates.

\paragraph{Thermal state and clock frequency.} A system-on-chip regulates its own temperature by
lowering its clock frequency when a threshold is reached, a behaviour called \emph{thermal
throttling}, and the board used here exposes a flag that records whether throttling
occurred~\cite{rpiosdocs}. A sustained inference
workload on all cores raises the temperature of a passively cooled board until throttling begins,
after which every timing figure measured on it is a figure for a slower clock. An active cooler
holds the clock at its rated frequency under the same load. A frequency governor, the
operating-system policy that chooses the clock in the absence of a thermal limit, can also lower
it to save power. A throughput figure measured on an \gls{sbc} is therefore a statement about the
board's thermal state and governor as much as about the model, and is interpretable only when
both are reported with it.

\paragraph{Throughput, latency and memory.} Three quantities describe the efficiency of a
configuration on the board. \emph{Throughput} is the number of output tokens decoded per second of
decode time. \emph{Latency} is the wall-clock time from the start of a request to a defined end
point, and for a language model it is split into the prefill latency, to the first output token,
and the decode latency, over the remaining tokens. Because latency varies between repetitions of
the same request with scheduling, cache state and thermal state, it is reported as percentiles of
a sample: the median, or p50, describes the typical case, and the 95th percentile, or p95,
bounds all but one repetition in twenty and is the figure a budget is set against. \emph{Peak
resident memory} is the largest amount of physical memory a process held during a run, and is the
figure a memory ceiling is checked against, since a model that fits on disk may still not fit
beside the other processes the board must run.

\clearpage
\section{Grammar-constrained decoding}
\label{sec:bg-gbnf}

A \emph{context-free grammar} is a set of rewriting rules over two alphabets: terminal symbols,
which appear in the strings of the language, and non-terminal symbols, each of which expands to a
sequence of terminals and non-terminals according to its rules. The language of the grammar is
the set of terminal strings derivable from its start symbol. \Gls{bnf} is the usual notation:
each rule names a non-terminal and lists its alternatives, separated by a bar. A grammar in which
no non-terminal can derive a string that contains itself, that is, a grammar without recursion,
has a finite language, and its longest string has a computable length. A \gls{json} object with a fixed set of
keys, enumerated string values and numbers of bounded width is such a language.

\paragraph{The token-masking mechanism.} Grammar-constrained generation restricts, at every
decoding step, the set of tokens the model is allowed to produce to those consistent with the
grammar, so that no output can leave the grammar's language regardless of what the model would
otherwise have produced. The runtime keeps the state of a parser over the output emitted so far;
before each step it determines which tokens of the vocabulary would extend that output to a
prefix of some string in the language, sets the logits of every other token to minus infinity,
and lets the sampler choose among the remainder. Greedy decoding then takes the most probable
\emph{admissible} token. Because a token is a piece of text that may span several characters,
and may straddle a boundary between two grammar symbols, the admissibility of a token is decided
character by character against the grammar. Formulations that compile a regular expression or a
context-free grammar into an automaton over the vocabulary make this check efficient, with
correctness results for the language classes the construction covers~\cite{willard2023,koo2024}.
\Gls{gbnf} is the \gls{bnf} dialect in which \texttt{llama.cpp} accepts a grammar: it adds
character classes and bounded repetition in the manner of regular expressions~\cite{llamacpp},
and is applied at every decoding step as described. The per-step cost
depends on the grammar's structure, and the runtime's documentation warns that some rule shapes
make sampling slow~\cite{llamacpp}.

\paragraph{Constraint against validation.} Two other ways of obtaining structured output exist,
and both differ from a grammar in kind. Prompting a model to produce well-formed output, and
checking the result afterwards, leaves the output's validity an empirical property: a rate
measured on some test set, which can change when the inputs change. Rejecting an invalid output
and decoding again converts that rate into a distribution of retries and latency. A grammar acts
inside the sampling loop, before an invalid token can be emitted rather than after a full
sequence has been produced, so validity holds on any input by construction, provided the length
cap exceeds the longest string in the language. What the grammar does not do is settle meaning. A
well-formed command addressed to the wrong aircraft is in the language, and no grammar can reject
it; the choice among admissible tokens remains the model's, and the masking also alters the
distribution the model samples from, since the probability mass of the excluded tokens is
redistributed among the admitted ones. Whether that alteration helps or harms accuracy on a given
task is a question for measurement, and Chapter~\ref{chap:state-of-the-art} reviews what has been
measured.

\section{Spoken-language understanding: intents and slots}
\label{sec:bg-slu}

A spoken command reaches a parser through a pipeline. A microphone captures audio; an endpointer
decides when an utterance has ended; a \gls{stt} model transcribes it; the parser maps the
transcript to a structured command; and a validator and controller act on the command. The
\gls{stt} stage in this work is a model of the Whisper family, an encoder--decoder transformer
trained on a large corpus of weakly supervised audio-transcript pairs~\cite{whisper}, served by a
C/C++ port of the same kind as the language-model runtime~\cite{whispercpp}. The pipeline's other
components, and
the audio chain around them, are the subject of the \emph{M\'emoire d'Ing\'enieur}; this
document concerns the parser, whose input is a transcript and whose output is a command.

\paragraph{The semantic frame.} Spoken-language understanding treats the meaning of a
task-oriented utterance as a \emph{semantic frame}: an \emph{intent}, drawn from a fixed set of
actions the system can perform, and a set of \emph{slots}, each a key with a value extracted
from the utterance~\cite{qin2021}. ``Form a circle with radius five metres'' has the intent
\emph{formation} and the slots \emph{shape} = circle and \emph{radius} = 5. Intent detection is
a classification problem over the utterance; slot filling is traditionally a sequence-labelling
problem over its words, and the field's methods are classified by whether the two are modelled
separately or jointly~\cite{qin2021}. Benchmark corpora define the frame the same
way: MASSIVE, for instance, labels a million utterances in 51 languages with one of 60 intents and
spans over 55 slot types~\cite{massive}. A generative formulation replaces the two classifiers by
one language model that emits the frame directly as structured text, for instance as a
\gls{json} object whose first key is the intent and whose remaining keys are the slots. Under
that formulation a grammar of the kind described in Section~\ref{sec:bg-gbnf} can fix which
intents and keys exist and what shape their values take.

\paragraph{Out-of-domain input and abstention.} A system that acts on commands must also handle
utterances that are not commands: speech addressed to someone else, a request outside the
vocabulary, or a transcript garbled by noise. The frame therefore includes a null intent,
conventionally \emph{unknown}, whose selection is a decision not to act. Whether a parser
selects it when it should is a property distinct from its accuracy on in-domain input, and the
two are scored separately (Section~\ref{sec:bg-metrics}). The cost of the two kinds of error is
asymmetric for a robot: a command declined can be repeated, whereas a command invented from
noise moves a vehicle.

\clearpage
\section{Evaluation metrics and paired comparison}
\label{sec:bg-metrics}

The metrics below are the standard instruments for the task and the setting described above.
Their operational definitions in this work, including the canonical form to which every output
is reduced before comparison, are fixed in Section~\ref{sec:definitions-of-record}.

\paragraph{Accuracy of a structured output.} \Gls{em} is the fraction of test items on which the
predicted frame equals the reference frame in its entirety, intent and every slot; one wrong slot
value makes the item wrong. It is the strictest of the metrics and the one a controller depends
on. Intent accuracy is the fraction of items whose intent alone is correct. Per-class precision,
the fraction of predictions of a class that are correct, and recall, the fraction of reference
items of that class that are found, combine into the F1 score, their harmonic mean; \emph{macro}
F1 averages the per-class F1 over the classes with equal weight, so that a frequent class cannot
dominate it, whereas \emph{micro} F1 pools the counts across classes. Slot F1 applies the same
precision and recall to the multiset of (key, value) pairs across all items. Schema validity is
the fraction of outputs that parse as a well-formed frame at all, independently of whether the
frame is correct; under a grammar it is unity by construction, and it is informative only where
the grammar is removed.

\paragraph{Safety-oriented rates.} Two rates score the abstention behaviour of
Section~\ref{sec:bg-slu}. The \emph{false-command rate} is the fraction of out-of-domain inputs
on which the system does anything other than decline; a fallback to a hold counts, since the
vehicle still acts on input that was not a command. The \emph{safe-failure
rate} is the fraction of the system's errors, on in-domain and out-of-domain input alike, that
resolve to a harmless action, such as declining or holding position, rather than to a wrong
command. The first has out-of-domain items as its denominator, the second erroneous items only,
whatever their input; neither is a function of \gls{em}, and a system can score well on \gls{em}
and poorly on both.

\paragraph{Speech recognition.} The accuracy of a transcript is scored by \gls{wer}: the number of
word substitutions, deletions and insertions needed to turn the transcript into the reference,
summed over a corpus and divided by the total number of reference words. An interval for a
corpus-level rate such as \gls{wer} is obtained by the bootstrap: the corpus is resampled with
replacement many times, the rate is recomputed on each resample, and the interval is read from
the percentiles of those rates. A \gls{wer} carries into the parser as noise in its input, so a
parser scored on recognised text is scored on the recogniser as well; which input each metric of
this document takes is fixed with its definition of record.

\paragraph{Paired comparison of two systems.} When two configurations are scored on the same
items, and each item's outcome is right or wrong, the two sets of outcomes are paired, and a test
that uses the pairing is more sensitive than one that compares the two accuracies as if they were
independent samples. McNemar's test does so: it discards the items on which both systems agree
and asks whether the \emph{discordant} items, those one system answers correctly and the other does not,
split evenly between the two directions~\cite{mcnemar1947}. With few discordant items the exact
binomial form of the test is used; with many, a chi-square approximation with a continuity
correction~\cite{edwards1948}. Two properties bound what the test can say. A non-significant
result is not evidence that the two systems are equivalent, only that the observed disagreement
is small or evenly split; with one discordant item the test cannot reject at all. And when
several comparisons are made from the same data, the chance that at least one reaches a nominal
significance level by accident rises with their number, which the Bonferroni correction controls
by dividing the significance level by the number of comparisons in the family~\cite{dunn1961}.

\paragraph{Trade-offs and dominance.} A comparison that scores configurations on more than one
criterion, accuracy and latency among them, need not produce a single winner. Configuration $a$
\emph{dominates} configuration $b$ when $a$ is at least as good as $b$ on every criterion and
strictly better on at least one. The configurations that no other dominates form the
\emph{Pareto frontier}, and a choice among them is a choice of how much of one criterion to give
up for the other. A selection rule that fixes a bound on one criterion, such as a latency budget,
and maximises the other within it, picks a point of that frontier, provided a tie on the
maximised criterion is broken by the bounded one; stated before the results exist, it keeps the
choice from being made to fit them.

\section{Conclusion}
\label{sec:bg-conclusion}

This chapter has set out the object of the comparison: a decoder-only \gls{slm} adapted by
low-rank fine-tuning. It has set out the three conditions under which it is compared: weight
quantisation to a \gls{gguf} format, decoding on the \gls{cpu} of an \gls{sbc} where memory
bandwidth and thermal state bound throughput, and token masking under a bounded \gls{gbnf}
grammar. It has defined the task, the extraction of an intent and its slots from a transcribed
command, with abstention as a required behaviour. And it has defined the instruments, from
\gls{em} and the F1 scores to McNemar's paired test and Pareto dominance. Chapter~\ref{chap:state-of-the-art} reviews what prior studies have measured
under each of these conditions, compares them to one another, and identifies what none of them
measures together.
